\section{Introduction}
A precise observation can reveal which signal to recover while
leaving a fixed iterative solver unable to reach it.
For two unit-separated templates along an eigenvector of $A^\top A$
with eigenvalue $\alpha$, their noiseless backprojections differ by
$\alpha$, whereas their observation separation in noise standard
deviations is $(\tau^2+\nu^2/\alpha)^{-1/2}$.
We study how denoising-scale choices interact with this gap in RED
gradient descent~\cite{romano2017}.

Graduated optimization has long used smoothing to guide visual
reconstruction~\cite{blake1987,mobahi2015,hazan2016}.
Fernsel et al.~\cite{fernsel2024} prove terminal stationarity for an
adaptive smoothing and line-search method and measure how initialization
and initial smoothing affect attainment of the global minimum in a
two-dimensional mixture. Hurault~\cite{hurault2026} proves terminal
stationarity for annealed RED--GD with nonincreasing scales and a
global Lipschitz step bound. Our comparison concerns recovery error of
the actual iterates across an entire class of admissible controls,
including arbitrary switching, and a finite-budget upper bound
after granting access to an intermediate band.

\paragraph{Can computation replace scale access?}
Theorem~\ref{thm:ind_noise} fixes the observation geometry, RED update,
and backprojection initialization. Across a positive-width box of
three-template models, every controller restricted to the two
specified variance ranges retains positive low-noise risk, uniformly
over all finite budgets. An observation-based preset using a fixed-width
intermediate band enters the correct terminal region in at most 25 calls
and then reaches vanishing risk. Corollary~\ref{cor:sensing_angle}
certifies the same separation for $|\varphi|\le.020$ radians at
$\alpha=10^{-3}$, uniformly as the noise vanishes.
Figure~\ref{fig:reachability} shows
the mechanism for the difficult observation class.

\paragraph{When do endpoints suffice?}
For two templates aligned with a sensing eigenvector, the endpoint
maps attain the maximum horizontal step before the midpoint.
An endpoint policy matches any strict crossing up to that midpoint
within the same budget.

\paragraph{How can a controller select useful scales?}
With three templates, prior mass and distance gaps predict local
drift and identify variances that can violate a candidate barrier.
Uniform channel and terminal checks then certify entry for all
positive-mass classes. We choose scales using these tests before
evaluating trajectories. Frozen photo-patch tests examine a
complementary case where different trajectories enter a common
terminal region and their differences decay.
\subsection{Related Work}
PnP and RED use denoisers to express priors for inverse
problems~\cite{venkatakrishnan2013,romano2017,kamilov2023}.
Convergence has been studied for PnP with continuation, Lipschitz
conditions, and MMSE denoisers~\cite{chan2017,ryu2019,xu2020}.
Generic RED residuals may lack a scalar potential~\cite{reehorst2019}.
The exact MMSE score identity~\cite{efron2011} and conditional
covariance calculus~\cite{dytso2023} give one here.
Near a terminal responsibility tie the RED Jacobian can grow as
$v_L^{-1}$, so our fixed step lies outside Hurault's sufficient
global step bound as $v_L\to0$.

Gaussian mixture scale space and mean shift describe mode formation
and tracking~\cite{carreira2003,carreira2007,carreira2015}.
Mixture priors also underpin patch restoration~\cite{zoran2011,yu2012}.
Leong and Traonmilin~\cite{leong2026} give schedule-dependent
recovery guarantees for varying projections under restricted
sensing and projection conditions. Our fixed RED update separates
statistical identifiability from reachability in a weak sensing
direction. Operator-linked PnP recovery~\cite{wolf2026} and
fixed-linear-denoiser contraction~\cite{sinha2025} provide
complementary perspectives. Annealed sampling, stochastic
regularization, and flow-based reconstruction use other updates
and objectives~\cite{song2019,kadkhodaie2021,laumont2022,renaud2024,martin2025}.

